\section{Ordinary normal tubes of a marked residue curve}
\label{sec:r13-cc-marked-curve-tubes}

Put \(R=K^\circ\), choose a uniformizer \(\pi\), and put
\(k=R/(\pi)\).  We retain the original field and its absolute value.
Let \(q\geq1\), and let \((Y,y)\) be a smooth integral model germ
of relative dimension \(q+1\), as in GE1.  Thus
\[
 \widetilde{\mathcal H(y)}=F,
 \qquad \operatorname{trdeg}_kF=q+1,
 \qquad |\mathcal H(y)^\times|=|K^\times|.
 \tag{CC1}
\]
Let \(V/k\) be a smooth integral model of \(F\), and let
\(C\subset V\) be a smooth proper integral curve.  Suppose that,
near the generic point \(\eta_C\), the curve is the transverse
intersection of labelled smooth prime divisors
\(D_1,\ldots,D_q\).  Choose their local equations
\(z_1,\ldots,z_q\in F\) at \(\eta_C\), and write
\(F_C=k(C)\).  Degrees of closed points of \(C\) are taken over
\(k\), including its full constant-field factor.

A finite norm-one central unit list \(a_1,\ldots,a_b\) is called
adapted to this curve and these normal parameters if
\[
 \bar a_j=z_1^{n_{1j}}\cdots z_q^{n_{qj}}u_j,
 \qquad n_{ij}\in\mathbf Z,
 \quad u_j\in\mathcal O_{V,\eta_C}^{\times}.
 \tag{CC2}
\]
Put \(g_j=u_j|_C\in F_C^\times\).  In particular every central
unit whose residue is a unit at \(\eta_C\) is adapted, with all
\(n_{ij}=0\).  This includes the simultaneous lifts supplied by ML.

\subsection{The exact zero slice}

Choose \(t\in\mathcal O_{V,\eta_C}^{\times}\) whose restriction
to \(C\) is a separating transcendence parameter.  The conormal
sequence makes
\(dz_1,\ldots,dz_q,dt\) a basis at \(\eta_C\).  A standard
etale neighborhood consequently has the form
\[
 B_0=\left(k[Z_1,\ldots,Z_q,T,G_0^{-1}][\Theta]/Q_0\right)_{H_0},
 \qquad \operatorname{Frac}(B_0)=F.
 \tag{CC3}
\]
The polynomial \(Q_0\) is monic, its derivative is invertible on
the selected chart, and all localizers are nonzero at \(\eta_C\).
After further localization in \(T\), its zero slice has a selected
coprime factorization
\[
 Q_0(0,T,\Theta)=p_0(T,\Theta)r_0(T,\Theta),
 \qquad k(T)[\Theta]/p_0=F_C.
 \tag{CC4}
\]
Invert the complementary factor on the selected component.  Lift
the coefficients to \(R\), choosing the lift so that
\[
                    Q(0,T,\Theta)=P(T,\Theta)R_0(T,\Theta)
                    \quad\hbox{exactly}.
 \tag{CC5}
\]
The modification used in DT4 works with the tuple \(Z\) in place
of its single variable: subtract the original zero specialization
and add the prescribed product.  This preserves the reduction and
monicity.

Let \(\mathscr W\) be the selected lifted smooth model defined by
\(Q=0\).  Let \(\mathscr C_0\) be the selected smooth affine
\(R\)-curve defined by \(P=0\), with all zero-slice localizers
inverted.  Write \(C_0=\mathscr C_{0,K}^{\rm an}\), and let
\(v\) be its generic special-fiber point.  Then
\[
 \widetilde{\mathcal H(v)}=F_C,
 \qquad |\mathcal H(v)^\times|=|K^\times|.
 \tag{CC6}
\]
The point \(v\) is considered in the boundaryless analytic curve
\(C_0\), with its ordinary neighborhoods.  Its proper residue
curve is the given \(C\), since both have function field \(F_C\).

\begin{theorem}[Ordinary normal product near every positive radius block]
\label{thm:r13-cc-ordinary-product}
There is an original ordinary comparison between the central
generic point of \(\mathscr W\) and \((Y,y)\), inducing the
specified residue-field identity.  The coordinates \(Z_i\) become
actual central units with residues \(z_i\).

In the selected chart there is a continuous family
\[
 x_C(r_1,\ldots,r_q),\qquad 0<r_i<1,
 \tag{CC7}
\]
identified, near every compact block
\(B\Subset(0,1)^q\), with the relative multivariable Gauss family
over \(v\) in \(\mathbf G_m^{q,\rm an}\times_K C_0\).
The identification is an isomorphism between ordinary open
subdomains and preserves \(Z_1,\ldots,Z_q,T\).
As all \(r_i\to1^-\), the family converges ordinarily to the
central point, uniformly with respect to the ratios of
\(-\log r_i\).  Hence a sufficiently short full collar lies
in any prescribed central ordinary comparison neighborhood.
\end{theorem}

\begin{proof}
The central comparison is the dimension-independent argument of
GE6.  Lift the finite residue generators on the original model,
solve the simple root by the ordinary analytic root lemma, and
compare the two etale parameter maps.  Their central completed
fields have the same residue field and the same discrete value
group as \(K\).  Successive \(\pi\)-adic residue approximation
makes the induced field inclusion surjective.  The pointed etale
germ theorem then gives an ordinary local isomorphism.

On the product side let \(\Theta_0\) be the zero-slice root.  At
each relative Gauss point over \(v\), (CC5) gives
\[
 |Q(Z,T,\Theta_0)|\leq\max_i r_i<1,
 \qquad |Q_\Theta(Z,T,\Theta_0)|=1.
 \tag{CC8}
\]
Every selected localizer and complementary factor has norm one.
On a compact block the first bound is uniformly less than one.
Ordinary continuity therefore gives a neighborhood on which the
Newton inequalities are uniform.  The retained analytic root
lemma gives the unique root \(\Theta\) in the prescribed strict
Hensel ball.  Solving \(P(T,\Theta_0)=0\) near \(\Theta\)
constructs the inverse.  Uniqueness proves the two compositions
are the identity after ordinary shrinking and glues the maps on
overlapping blocks.  This argument uses no independence condition
on the radii.

For convergence, complete the selected coordinate ring along the
normal ideal.  The simple root construction embeds it in the
formal series ring over the zero-slice coefficient algebra.  Its
reduction embeds the integral local chart into
\(F_C[[z_1,\ldots,z_q]]\).  If an integral function \(b\) has
nonzero reduction, some coefficient of its normal expansion has
nonzero residue in \(F_C\), and hence norm one at \(v\).
All coefficients have norm at most one.  For one such exponent
\(\alpha\),
\[
          r^\alpha\leq |b(x_C(r))|\leq1.
 \tag{CC9}
\]
After dividing a general nonzero integral function by its maximal
\(\pi\)-power, the same estimate applies.  Quotients give
convergence of every regular-function seminorm.  This proves
ordinary convergence, and the lower bound depends only on
\(\alpha\), not on ratios among the positive logarithmic radii.
\end{proof}

\subsection{The coefficient field at commensurable radii}

\begin{lemma}[The full Gauss residue extension]
\label{lem:r13-cc-gauss-residue}
For \(r\in\mathbf R_{>0}^q\), let \(K_r\) be the completion of
\(K(Z_1,\ldots,Z_q)\) for its Gauss norm.  Put
\[
 \Lambda_r=\{(a_0,a)\in\mathbf Z^{q+1}:
                     |\pi|^{a_0}r^a=1\}.
 \tag{CC10}
\]
If \(\ell_1,\ldots,\ell_d\) is a basis of \(\Lambda_r\),
the residues of the corresponding monomials
\(\pi^{\ell_{j0}}Z^{\ell'_j}\) are algebraically independent,
and
\[
 \widetilde {K_r}=k(U_1,\ldots,U_d),\qquad
 |K_r^\times|=|K^\times|r_1^{\mathbf Z}\cdots r_q^{\mathbf Z}.
 \tag{CC11}
\]
At \(x_C(r)\) one has
\[
 \widetilde {\mathcal H(x_C(r))}=F_C(U_1,\ldots,U_d),
 \qquad |\mathcal H(x_C(r))^\times|=|K_r^\times|.
 \tag{CC12}
\]
For a closed point \(p\in C\), its base change to
\(k(U_1,\ldots,U_d)\) remains one closed point of degree
\(\delta_p=[\kappa(p):k]\), and all orders of functions in
\(F_C^\times\) at that point are unchanged.
\end{lemma}

\begin{proof}
The associated graded field of a Gauss rational-function field
is the graded fraction field of
\(\operatorname{gr}K[\widetilde Z_1^{\pm1},\ldots,
\widetilde Z_q^{\pm1}]\).  A degree-one Laurent monomial has
exponent in \(\Lambda_r\).  Factoring one monomial from any
homogeneous numerator and denominator shows that the degree-one
subfield is exactly \(\operatorname{Frac}(k[\Lambda_r])\).
The exponent group is free, so this is the displayed rational
function field.  Distinct Laurent monomials in a basis have
distinct exponents, proving algebraic independence.  Completion
adds no initial forms: every nonzero completed element has an
approximation with strictly smaller error than its norm.

The identical calculation with the coefficient field
\(\mathcal H(v)\), whose residue and value group are (CC6),
gives (CC12).  The analytic product isomorphism transports that
calculation to the actual point.  In particular the relative
value-group index is one at every radius vector.

Finally \(\kappa(p)\otimes_k k(U)=\kappa(p)(U)\), a field of
degree \([\kappa(p):k]\).  Flat base change of the smooth curve
preserves its local uniformizer and the orders of functions
defined over \(k\).  The corresponding proper residue curve
over \(\widetilde {K_r}\) is \(C_{k(U)}\), with the displayed
function field.  This also explains the extra residue variables
at commensurable radii.
\end{proof}

\subsection{Finite central normalization and the complete local source}

Choose a finite tuple of rational functions on \(C\) which,
together with the \(g_j\), generates \(F_C\).  Lift each through
\(\mathcal O_{V,\eta_C}\twoheadrightarrow F_C\) to a unit,
and then to an actual central analytic unit by the retained
finite residue-unit lifting theorem.  Denote the resulting
enlarged adapted list again by \(a_1,\ldots,a_N\).
Choose actual rational units \(R_j\) near \(v\) on \(C_0\)
with residues \(g_j\).

\begin{theorem}[All-radii split tuple and complete marked block]
\label{thm:r13-cc-complete-marked-block}
There is \(r_*<1\) such that for each compact nonempty block
\(B\Subset(r_*,1)^q\) one ordinary product-comparison domain
contains the Gauss source over \(B\) times a neighborhood of
the complete curve source germ at \(v\), and on that domain
\[
 \left|\frac{a_j}{\prod_i Z_i^{n_{ij}}}-R_j\right|<|R_j|,
 \qquad
 t_{a_j}=\sum_i n_{ij}t_{Z_i}+t_{R_j}.
 \tag{CC13}
\]
The normalized graded initial forms agree there.

For the completed tuple
\(\mathcal A=(Z_1,\ldots,Z_q,a_1,\ldots,a_N)\), the entire
full-rank source in this domain is the relative Gauss product
of the complete curve source of \(\mathbf R\) with the radial
parameter domain.  Its lattice and calibration are the product
lattice and product calibration.  The graph (CC13) has integral
index one on every face and every quotient by the radial lattice.

Let
\[
 M=\bigcup_j\operatorname{Supp}\operatorname{div}_C(g_j),
 \quad b_p=(\operatorname{ord}_p(g_j))_j,
 \quad m_p=\gcd_j |\operatorname{ord}_p(g_j)|,
 \quad v_p=b_p/m_p.
 \tag{CC14}
\]
This is a finite marked set with at least two points, because the
residue tuple contains a nonconstant function.  At every core point over
\(B\), the complete labelled transverse normal one-cycle is
\[
              \sum_{p\in M}\delta_pm_p
                              [\mathbf R_{\geq0}v_p].
 \tag{CC15}
\]
Each label is an actual branch germ.  Repeated geometric ray
images are aggregated only after their labelled calibrations.
For every normalized coordinate,
\[
       w_p(v_p)_j=\delta_p\operatorname{ord}_p(g_j),
       \qquad w_p=\delta_pm_p.
 \tag{CC16}
\]
Every branch of the completed tuple is retained, including those
whose projection to an earlier tuple is zero or rank deficient.
\end{theorem}

\begin{proof}
Represent \(u_j\) by a quotient regular and nonzero at
\(\eta_C\).  In the lifted chart its specialization to the zero
slice is a curve unit with residue \(g_j\).  Its difference from
that specialization belongs to the normal ideal in the formal
expansion; it has norm less than one at every positive-radius
core point.  Any other chosen lift \(R_j\) differs by a
zero-residue function at \(v\).  Finally the original central
unit and its chosen rational residue representative have ratio
with residue one at \(y\), so their difference has smaller norm
on an ordinary neighborhood of \(y\).  Convergence (CC9) puts
one full collar for the finite list inside this neighborhood.
These facts prove the strict inequality at every core point, including commensurable radii.  Continuity and compactness
extend it to one ordinary domain over \(B\).

For the split tuple \((Z,\mathbf R)\), the retained TP graded-rank
argument applies with \(q\) radial variables and a one-dimensional
factor.  Indeed full graded rank \(q+1\) forces rank one in the
curve coordinates and graded independence of all \(q\) radial
coordinates over the whole completed coefficient field.  This
is exactly the relative Gauss criterion.  Conversely a relative
Gauss point over a rank-one curve source has the required rank.
This identifies the entire source, at every radius vector.
The finite graph (CC13) has a coordinate projection as an
integral retraction, so it changes neither source rank nor any
of the stated lattice indices.  The product calibration is the
actual curve calibration times the primitive torus calibration
one.  The relative Gauss parameterization is continuous by the
maximum formula.  Applying the topological tube lemma to its
source parameter space over the compact set \(B\) gives a common
curve-source neighborhood.  No Cartesian basis for arbitrary
ambient ordinary neighborhoods is used.

The curve source at the interior point \(v\) is the complete
proper-residue-curve source of the retained AC/FP/NCB theorems.
Its visible branches are exactly the finite divisor support in
(CC14), and there is one actual source halfedge in each such
branch.  Its branch degree is \(\delta_p\), its coarse index is
one, and its primitive stretching is \(m_p\).  Thus its weight
is \(\delta_pm_p\).  The product-source identity gives the same
labelled normal block at every point over \(B\).  Alternatively,
(CC11)--(CC12) compute directly the curve field after a radial
base change: the marked points keep degree \(\delta_p\), their
orders are unchanged, and the relative value-group index remains
one.  No residue-field equality with \(F_C\) is asserted when
\(d>0\) in (CC10).

Finally the radial columns are
\(l_i=(e_i,n_{i1},\ldots,n_{iN})\).  Projection to the first
\(q\) coordinates splits their lattice integrally.  Its quotient
is identified with the normalized curve-coordinate lattice by
subtracting \(\sum_i n_{ij}t_{Z_i}\).  This proves the primitive
normal formula (CC16) and retains every completed branch.
\end{proof}

\begin{corollary}[Finite adapted enlargements]
\label{cor:r13-cc-adapted-enlargement}
For any further fixed finite adapted central unit list, the theorem
holds on a common collar and common finite source subdivision.
Its marked set is the union of the old divisor supports with the
new ones.  At an old mark, if \(m'_p\) is the new gcd, then
\(m'_p\mid m_p\); forgetting the new coordinates maps the new
primitive ray with index \(m_p/m'_p\), and
\[
              (m_p/m'_p)\,\delta_pm'_p=\delta_pm_p.
 \tag{CC17}
\]
A new mark has zero old transverse order vector and remains an
actual branch of the enlarged block.
\end{corollary}

\begin{proof}
Apply (CC13) to the union of the two finite lists and intersect
their ordinary domains.  The residue order vectors are enlarged
by adjoining rows, and their projection is exactly the old order
vector.  Dividing by the respective gcd gives (CC17).  The
canonical finite local source structure supplies a common
subdivision for these finitely many unit-log functions.
\end{proof}

\subsection{The old tangent index}

\begin{proposition}[Factorization of the source projection index]
\label{prop:r13-cc-tangent-index}
Suppose projection to an earlier moment lattice maps the completed
radial lattice \(\mathbf Z^q\) injectively to \(A\mathbf Z^q\)
in a saturated rank-\(q\) lattice \(N_T\).  Put
\(i_T=[N_T:A\mathbf Z^q]\).  On a transverse branch with nonzero
old quotient image, let \(m_p^{\rm old}\) be the primitive
stretching of its integral order vector in the saturated old
quotient lattice.  Then the full source projection index on that
branch is
\[
                       i_T\,m_p^{\rm old}/m_p,
 \tag{CC18}
\]
and the projected calibration is
\(\delta_p i_Tm_p^{\rm old}\).
\end{proposition}

\begin{proof}
Let \(L\) be the lattice generated by the projected radial
columns and the projected completed primitive transverse vector.
Its intersection with \(N_T\) is \(A\mathbf Z^q\), since the
transverse quotient vector is nonzero.  The exact sequences of
\(L\) and of its saturation, with their tangent sublattices,
give the product of the tangent index \(i_T\) and the quotient
index \(m_p^{\rm old}/m_p\).  The latter is an integer because
the projection is integral.  Multiply by the completed weight
\(\delta_pm_p\), and apply the retained calibrated finite moment
projection formula.  This is a source lattice calculation; no
proper pushforward of an arbitrary current is involved.
\end{proof}

\subsection{Ordinary local constancy of algebraic curve residues}

\begin{lemma}[All-radius polynomial constancy]
\label{lem:r13-cc-polynomial-constancy}
Let \(a\) be a norm-one actual central unit whose residue is a
unit at \(\eta_C\), with nonzero restriction \(g\in F_C\)
algebraic over \(k\).  On one sufficiently short collar, every
point \(x_C(r)\) has an original ordinary neighborhood on which
\(|a|=1\).  A finite list admits the same collar and simultaneous
ordinary neighborhoods.  The assertion holds for every radius
vector, including all commensurable ones.
\end{lemma}

\begin{proof}
Let \(H\in R[T]\) be a monic coefficient lift of the minimal
polynomial of \(g/k\).  Its constant term is a unit.  On the
zero-slice curve, \(H(R_g)\) has zero residue at \(v\), hence
norm less than one.  The strict comparison (CC13) gives
\(|H(a)|<1\) at every core point in one finite collar.  By
continuity this inequality holds on an ordinary neighborhood.
If \(|a|>1\), the monic term of \(H(a)\) strictly dominates all
other terms; if \(|a|<1\), the unit constant term dominates.
Both alternatives contradict \(|H(a)|<1\).  Thus \(|a|=1\).
Finite intersections prove the simultaneous version.
\end{proof}

\begin{corollary}[An ordinary image in a supplied tangent plane]
\label{cor:r13-cc-tangent-plane-constancy}
Suppose the order columns of an old tuple span a real plane
\(L\), and each character in an integral basis of its annihilator
has zero normal orders and algebraic nonzero restriction to
\(C\).  Then every sufficiently near core point has an ordinary
neighborhood whose old moment image is contained in \(L\).
At a point in the relative interior of a specified open cone in
\(L\), the neighborhood may be chosen with image in that cone.
\end{corollary}

\begin{proof}
Apply the preceding lemma to the finitely many actual monomial
characters.  Their logarithms vanish identically on a common
ordinary neighborhood; their common real kernel is \(L\).
The last assertion follows by ordinary continuity.
\end{proof}

\paragraph{Source dependencies.}
The construction uses the reviewed DT/GE ordinary analytic root
and central comparison arguments, TP's complete relative Gauss
source and its calibration, and the original curve branch-degree
and finite-support inputs AC/FP as retained in NCB.  The new
features are the simultaneous normal zero slice, the explicit
commensurable-radius residue computation, and the complete
source-localized curve block for any supplied normal rank \(q\).
The finite Picard completion and the comparison with other actual
source blocks are recorded separately.
